% part: intuitionistic-logic
% chapter: semantics
% section: introduction

\documentclass[../../../include/open-logic-section]{subfiles}

\begin{document}

\olfileid{int}{sem}{int}

\olsection{جان پہچان}

معنیات توں بناں کسے منطق دا وی تسلی بخش بیان نہیں ہو سکدا۔ ایہ گل وجدانیت پسند
منطق اُتے وی لاگو ہوندی اے۔ کلاسیکی منطق وچ !!{valuation}s اُتے ٹکی معنیات
معیاری اے، پر وجدانیت پسند منطق لئی کئی وکھ وکھ معنیات پیش کیتیاں گئیاں نیں۔
ایہناں وچوں کوئی وی ایس لحاظ نال پوری طرح تسلی بخش نہیں کہ اوہ منطقی رابطیاں
دے معنی وجدانیت پسند منطق دے اصولاں مطابق قبول ہون جوگے طریقے نال سمجھاوے۔

!!{relational model}s اُتے ٹکی معنیات، جیہڑی امکانی منطق دی معنیات ورگی اے،
شاید سبھ توں ودھ رائج اے۔ ایس معنیات وچ !!{propositional variable}s نوں جہاناں
نال جوڑیا جاندا اے تے جہان آپس وچ اک رسائی تعلق نال جڑے ہوندے نیں۔
اوہ تعلق ہمیشہ جزوی ترتیب ہوندا اے، یعنی انعکاسی، ضدِ تقابلی تے متعدی۔

سمجھن لئی تسی ایہناں جہاناں نوں علم دیاں حالتاں یا ’گواہی دیاں صورتاں‘ سمجھ سکدے او۔
حالت~$w'$ تک~$w$ توں رسائی اُدوں تے صرف اُدوں اے جدوں ساڈے موجودہ علم دے لحاظ
نال $w'$ علم دی اک ممکنہ (اگلی) حالت ہووے، یعنی~$w$ وچ جو کجھ معلوم اے
اوہدے نال میل کھاوے۔ اک واری کوئی قضیہ معلوم ہو جاوے تاں اوہ مُڑ نامعلوم نہیں
ہو سکدا۔ یعنی جدوں وی $!A$ دا~$w$ وچ پتا ہووے تے~$Rww'$ ہووے، تاں $!A$ دا
~$w'$ وچ وی پتا اے۔ سو رسائی تعلق دے لحاظ نال ’علم‘ گھٹدا نہیں۔

جے اسیں علم دی منطق وانگ ’$!A$ معلوم اے‘ دی تعریف ’سارے علم دے متبادلاں وچ سچ اے‘
دے طور اُتے کریے، تاں $!A \land !B$،~$w$ وچ معلوم اے جے سارے علم دے متبادلاں وچ
$!A$ تے~$!B$ دونیں معلوم نیں۔ پر علم گھٹدا نہیں تے $R$ انعکاسی اے، ایس لئی
$!A \land !B$، $w$ وچ اُدوں تے صرف اُدوں معلوم اے جدوں $!A$ تے $!B$،~$w$
وچ معلوم نیں۔ ایس وجہ نال $!A \lor !B$ وی $w$ وچ اُدوں تے صرف اُدوں معلوم
اے جدوں دوناں وچوں گھٹ توں گھٹ اک معلوم ہووے۔ سو $\land$ تے $\lor$ دے سچ ہون
دیاں شرطاں اوہو نیں جیہڑیاں کلاسیکی منطق وچ نیں۔

پر شرطیے دے سچ ہون دیاں شرطاں کلاسیکی منطق توں وکھریاں نیں۔ $!A \lif !B$،
~$w$ وچ اُدوں تے صرف اُدوں معلوم اے جدوں کوئی $w'$، جیہڑی $Rww'$ پوری کرے، ایہو جیہی
نہ ہووے جتھے $!A$ معلوم ہووے تے $!B$ معلوم نہ ہووے۔ ایہ ایس شرط ورگی نہیں
کہ $!A$ نامعلوم ہووے یا $!B$ معلوم ہووے، ایہ گل~$w$ وچ ہووے۔ کیوں جے، جے سانوں نہ $!A$ دا پتا اے نہ $!B$ دا، ایہ گل~$w$ وچ ہووے، تاں ہو سکدا اے اگلی علمی حالت~$w'$، جیہڑی $Rww'$ پوری کرے، ایہو جیہی ہووے کہ $w'$ وچ $!A$ معلوم ہو جاوے پر~$!B$ دا پتا نہ لگے۔

سانوں $\lnot !A$ دا پتا صرف اوہدوں اے جدوں اگوں علم دی کوئی ممکنہ حالت
ایہو جیہی نہ ہووے جتھے سانوں~$!A$ دا پتا ہووے۔ گل ایہ اے کہ جے $!A$ دا
پتا لگ سکدا ہووے، تاں علم دی کسے ممکنہ اگلی حالت وچ~$!A$ معلوم ہو جاوے گا۔
چونکہ سانوں $\lfalse$ دا پتا نہیں لگ سکدا، ایس لئی اوہ اگلی علمی حالت وچ
سانوں $!A$ دا پتا ہووے گا پر~$\lfalse$ دا نہیں۔

ایس تعبیر وچ تیسرے امکان دے اخراج دا اصول نہیں چلدا۔ کیوں جے کجھ $!A$ ایسے
ہوندے نیں جیہناں دا سانوں ہجے پتا نہیں، پر پتا لگ سکدا اے۔ ایہو جیہے
!!a{formula}~$!A$ لئی $!A$ تے~$\lnot !A$ دونیں نامعلوم نیں، ایس لئی
$!A \lor \lnot !A$ وی نامعلوم اے۔ پر مثال لئی ایہ گل سانوں معلوم اے کہ
$\lnot(!A \land \lnot !A)$۔ کیوں جے کوئی اگلی حالت ممکن نہیں جتھے سانوں
$!A$ تے $\lnot !A$ دوناں دا پتا ہووے، تے ایہ گل سانوں معلوم اے بھانویں سانوں
~$!A$ یا $\lnot !A$ دا پتا ہووے بھانویں نہ ہووے۔

!!^{relational model}s توں علاوہ وجدانیت پسند منطق دی ہور معنیات وی اے۔
اک ہور ٹوپولوجی والی معنیات اے۔ ایتھے قضیاں دی تعبیر ٹوپولوجی والی فضا دے
کھلے سیٹاں دے طور اُتے ہوندی اے تے رابطیاں دی تعبیر ایہناں سیٹاں اُتے عملاں دے
طور اُتے (مثلاً $\land$ دا میل اشتراک نال اے)۔

\end{document}
